\subsection{Арифметические операторы}

Python поддерживает все стандартные математические операции.

\begin{itemize}
    \item \texttt{+} Сложение
    \item \texttt{-} Вычитание
    \item \texttt{*} Умножение
    \item \texttt{/} Деление (всегда возвращает \texttt{float})
    \item \texttt{//} Целочисленное деление (отбрасывает дробную часть)
    \item \texttt{\%} Остаток от деления
    \item \texttt{**} Возведение в степень
\end{itemize}

\begin{codefile}[python]{python/01_python_basics/03_arithmetic.py}
\end{codefile}

Обратите внимание на приоритет операций: умножение и деление выполняются раньше сложения и вычитания, как в математике. Скобки \texttt{()} используются для изменения порядка действий.

\begin{taskbox}[Калькулятор]
Напишите программу, которая считает площадь прямоугольника со сторонами 5 и 8 метров.
\end{taskbox}

\subsubsection*{Задачи для самостоятельного решения}

\begin{enumerate}
    \item Сложите числа 15 и 38. Выведите результат.
    \item Вычтите из 100 число 45. Выведите результат.
    \item Умножьте 7 на 8.
    \item Разделите 50 на 4. Какой тип данных получился?
    \item Сколько раз число 4 помещается в 50 целиком? Используйте \texttt{//}.
    \item Чему равен остаток от деления 20 на 3? Используйте \texttt{\%}.
    \item Возведите 2 в степень 10 (2 в 10-й степени).
    \item Вычислите \texttt{2 + 2 * 2}. Затем \texttt{(2 + 2) * 2}. Сравните результаты.
    \item Сторона квадрата \texttt{a = 5}. Найдите периметр \texttt{P = 4 * a}.
    \item Даны стороны \texttt{a = 10}, \texttt{b = 20}. Найдите площадь \texttt{S = a * b}.
    \item Сколько секунд в одном часе? (60 секунд * 60 минут).
    \item Даны числа 10, 20, 30. Найдите их среднее арифметическое.
    \item Дано число 12345. Найдите его последнюю цифру, используя оператор \texttt{\% 10}.
    \item Дано четное число \texttt{n = 10}. Получите следующее четное число (+2).
    \item Умножьте -5 на -5. Проверьте знак результата.
\end{enumerate}
